\section{How to Break an Egg}

\begin{quotex}
The person who writes for fools is always assured of a large audience.
\flright{\textsc{Arthur Schopenhauer}}
\end{quotex}

Lilliput is a small island nation in the Indian Ocean. They have been wracked for centuries by an unresolvable dilemma: should you break a soft-boiled egg on the big end or the little end? The traditional practice was to break them on the larger end until a radical emperor decreed that they should be broken on the small end. This led to a series of internal conflicts that continue to this day. At first, scientists tried to answer that question. Then theologians scanned their scriptures looking for answers. Personal preference was not allowed by law. One side had tradition in its favor, the other did not.

Now that social media has reached the island, the battle is played out in public. That led to such bitterness that two competing social media platforms have arisen: one for big-enders and the other for little-enders. Now they don't even need to talk to each other.

\paragraph{Vaccine destroys the image of God}


A renegade Orthodox priest in a podcast recently concluded that Covid vaccines alter one's DNA, thereby destroying the image of God in the vaccinated. The notion that God favors certain DNA is hardly new, but not so blatantly stated in recent years.

\paragraph{High Culture}


Human beings are meant to live in a civilization. So-called primitive people are not earlier versions of the human race, but rather the degenerate remnants of a higher civilization. Otherwise, the human race could not have survived in the wild.

\paragraph{The Chinese Imperial Examination}


The Traditional view is that the social system should be guided --- but not ruled --- by the spiritual class. That served to maintain ancient civilizations like China, India, and old Egypt for centuries. India has had a hereditary caste of Brahmins to fill that role. China, on the other hand, relied on an examination system from about the seventh to the early twentieth century. Thus, people from all classes were able to compete.

The exam was so strict that only one in a thousand could pass. In principle, the notion that the best and the brightest should rule sounds fine. There are several high IQ societies that accept only the top 99.9 percentiles. However, intelligence alone is not enough, so that these societies have accomplished nothing of merit. Currently in the West, the spiritual guides are selected from certain elite universities. Unfortunately, they are not being trained to sustain the nation's culture but rather to overturn it.

The Chinese imperial examination\footnote{\url{https://en.wikipedia.org/wiki/Imperial\_examination}}, on the other hand, included questions on Confucian philosophy, ethics, and culture. That served to protect the spiritual foundations of the Chinese people.

\paragraph{Natural Law}


Western ethics used to be based on the concept of Natural Law, viz., one should act in conformance with one's true nature. For example, if man is a ``rational animal'', then he must act intelligently in all situations, not emotionally nor motivated by sensual desire. Moreover, anything that diminishes one's intellectual capacity, such as alcohol or drugs, is inherently immoral.

Thomas Aquinas developed a hierarchy of loyalties that is remarkably in conformance with biology and group selection, as understood today. Revolutionaries, including leading clerics in the Church, teach the opposite which is the foundation of revolutionary ideology.

\paragraph{Great Divergence}


The Great Divergence\footnote{\url{https://en.wikipedia.org/wiki/Great\_Divergence}} is the term used to describe the sudden dominance of Europe over the rest of the world, beginning in the Renaissance and continuing until recently. The elite view today is that the dominance was achieved through racism which hardly makes sense. On the contrary, it does not illustrate the superiority of the European races, since Asian nations have been achieving the same results.

The West seems to be undergoing the lassitude of Mouse Utopia.

\paragraph{Capital expenditures}


It used to be that a nation's power depended on capital formation (CapEx), such as factories, farms, mines, in short, anything that could produce wealth. Now, the West is converting to a system of Operational Expenditure (OpEx), in which capital is neither created nor preserved. Because of tax laws, it is often beneficial to rent rather than buy. This is happening quickly today in cloud technology, in which a few companies --- Amazon, Microsoft, and Google, primarily --- are investing capital in computer technology, while other companies are disbanding their computer infrastructure and renting computer power from the cloud companies.

Trump's tax cut a few years ago did the opposite of what was intended. After hearing an economist on TV praised the plan, I commented on his Twitter feed that companies would use the windfall to buy back stocks (a practice that used to be illegal) rather than create capital. He ignored me.

Now we see the results today with the worldwide chip shortage. Intel, which used to make its own chips in the USA, spent \$100 billion gained from the tax cut to buy back stock. That raised the price for investors. They should have used that money to build a new chip manufacturing facility.

Capital accumulation is an advantage to a country, but Marxists claim it leads to inequality. Just so you know how the tell the players apart.

\paragraph{Common Task}


The oddball 19\textsuperscript{th} century Russian philosopher, Nikolai Fedorov, developed the idea of the Common Task. This is how he defined it:

\begin{quotex}
There is not a single invention which the military are not bent on applying to warfare, not a single discovery which they fail to turn to military purposes. So if it were made the duty of the armies to adapt everything now used in warfare for controlling natural forces, this duty would automatically become the common task of humanity as a whole.
\end{quotex}


His point is just as valid today. Theological problems may be better served by technology and administration. For example, once you get past your food snobbery, McDonalds has done more to ``feed the poor'' than the church or government. The same goes for ``healing the sick''; for the government to promise ``healthcare for all'', there has to be healthcare in the first place.

Fedorov is in the unique position as being a strong influence both on the Russian Sophiologists as well as on transhumanists.

\paragraph{Hunter gatherers}


The Platonic ideal of justice is ``to each his own''. The hunter-gatherer ideal is ``see and take''.

\paragraph{Sex and iguanas}


The iguanas have been skittish lately, seldom coming out in public. However, it must be rutting season, because, in their excitement, they are mating in broad daylight, making themselves vulnerable. My son was able to get a kill on the male but the female scampered away. I asked him if she asked for a cigarette before she left.

\url{https://www.youtube.com/watch?v=nODFXvTUyKM}

\flrightit{Posted on 2022-01-04 by Cologero}

